\documentclass[11pt, aspectratio=169]{beamer}

% ==========================================
% Theme & Color Settings (Moloch Theme)
% ==========================================
\usetheme{moloch}
\usecolortheme{moloch}

% ==========================================
% Modular Preambles & Configurations (Relative)
% ==========================================
% ==========================================
% preamble/packages.tex
% Essential Packages for Thesis Beamer Slides
% ==========================================

% Math Extensions
\usepackage{amsmath}
\usepackage{amssymb}
\usepackage{amsfonts}

% Tables Formatting
\usepackage{booktabs}
\usepackage{array}

% Code and Verbatim Environments
\usepackage{listings}

% Vector Graphics (TikZ)
\usepackage{tikz}
\usetikzlibrary{calc, positioning, shapes.geometric, arrows.meta}

% Miscellaneous utilities
\usepackage{graphicx}
\usepackage{hyperref}

% ==========================================
% preamble/macros.tex
% Custom Math Operators and Thesis Shortcuts
% ==========================================

% Math symbols and operators
\newcommand{\E}{\mathbb{E}}
\newcommand{\R}{\mathbb{R}}

% Vector notation shortcut
\newcommand{\vect}[1]{\mathbf{#1}}

% Algorithmic symbols
\newcommand{\norm}[1]{\left\lVert#1\right\rVert}
\newcommand{\grad}{\nabla}

% Slide emphasis helpers
\newcommand{\keyterm}[1]{\textbf{\alert{#1}}}

% ==========================================
% preamble/listings-setup.tex
% Syntax Highlighting for Code Listings
% ==========================================

% Curated modern colors for listings
\definecolor{codegray}{rgb}{0.5,0.5,0.5}
\definecolor{codeblue}{rgb}{0.13,0.29,0.53}
\definecolor{codegreen}{rgb}{0.18,0.55,0.34}
\definecolor{codeorange}{rgb}{0.85,0.45,0.10}
\definecolor{codebackground}{rgb}{0.97,0.97,0.98}

\lstset{
    backgroundcolor=\color{codebackground},
    basicstyle=\ttfamily\small,
    breakatwhitespace=false,
    breaklines=true,
    captionpos=b,
    commentstyle=\color{codegray}\itshape,
    keepspaces=true,
    keywordstyle=\color{codeblue}\bfseries,
    numberstyle=\tiny\color{codegray},
    stringstyle=\color{codeorange},
    showspaces=false,
    showstringspaces=false,
    showtabs=false,
    tabsize=4,
    frame=single,
    rulecolor=\color{codegray!30},
    frameround=tttt
}


% ==========================================
% Presentation Metadata
% ==========================================
\title{FedMAQ: Multi-Adaptive Quantization \& Knowledge Distillation}
\subtitle{Communication-Efficient Federated Learning}
\author{Christian Joseph Bunyi}
\institute{Dr. Andrew L. Tan Data Science Institute \\ De La Salle University}
\date{\today}

% ==========================================
% Slide Deck Document Structure
% ==========================================
\begin{document}

% Slide 1: Title Slide
\begin{frame}[plain]
    \titlepage
\end{frame}

% ==========================================
% Section 1: The Strategic Pivot
% ==========================================
\begin{frame}{Research Status}
    \begin{itemize}
        \item We shifted focus from physical IoT/Smart Campus energy forecasting to classical Federated Learning.
        \item We chose image classification (e.g., MNIST, CIFAR) to make experiments more manageable and reproducibility easier.
        \item We want to explore multi-adaptive quantization with server-side knowledge distillation with the hopes of further reducing communication while maintaining comparable accuracy.
    \end{itemize}
\end{frame}

% ==========================================
% Section 2: Background & Baseline Inspiration
% ==========================================
\begin{frame}{(Some) Related Works on Communication-Efficient FL}
    \centering
    \vfill
    \begin{table}
        \fontsize{7pt}{9.5pt}\selectfont
        \begin{tabular}{lcccc p{3.6cm}}
            \toprule
            \textbf{Algorithm}     & \textbf{Time-Aware} & \textbf{Data-Aware} & \textbf{Resource-Aware} & \textbf{Uses KD} & \textbf{Key Components}                 \\
            \midrule
            DynFed                 & Yes                 & Yes                 & No                      & Yes              & Quantization-Aware KD                   \\
            AdaDQ-KD               & No                  & No                  & Yes                     & Yes              & Dithering Quantization \& Diff. Privacy \\
            DAdaQuant              & Yes                 & Yes                 & No                      & No               & Doubly-Adaptive Fixed-Point             \\
            FedDT                  & No                  & No                  & No                      & Yes              & Ternary Compression                     \\
            LAQ-HC                 & No                  & Yes                 & Yes                     & No               & Reinforcement Learning                  \\
            FedPAQ                 & No                  & No                  & No                      & No               & Periodic Averaging                      \\
            \midrule
            \textbf{FedMAQ (Ours)} & \textbf{Yes}        & \textbf{Yes}        & \textbf{Yes}            & \textbf{Yes}     & \textbf{3D-Awareness + Logit KD}        \\
            \bottomrule
        \end{tabular}
    \end{table}
    \vspace{0.15cm}

    {\fontsize{8pt}{9.5pt}\selectfont We explore considering client resource constraints, data skew, and temporal training dynamics.}
    \vfill
\end{frame}

% ==========================================
% Section 3: The FedMAQ Architecture
% ==========================================
% \begin{frame}{FedMAQ Architecture}
%     \begin{columns}[T]
%         \begin{column}{0.48\textwidth}
%             \small
%             \begin{itemize}
%                 \item The server broadcasts the global model parameters to participating client nodes.
%                 \item Clients perform local optimization steps on their non-IID datasets.
%                 \item Uplinks are compressed using multi-adaptive quantization based on live bandwidth and stability.
%                 \item The server aggregates client logits via knowledge distillation to neutralize quantization noise.
%             \end{itemize}
%         \end{column}
%         \begin{column}{0.48\textwidth}
%             \centering
%             \includegraphics[width=\textwidth,height=0.65\textheight,keepaspectratio]{figures/dynfed_fig2.png}
%             \vspace{0.1cm}

%             {\fontsize{5.5pt}{7pt}\selectfont He et al. (2025 - DynFed), Figure 2: Overview of the training procedure using logits.\\}
%         \end{column}
%     \end{columns}
%     \vfill
%     \centering
%     {\fontsize{8pt}{9.5pt}\selectfont Our work is a logical extension of DynFed by He et al (2025).}
% \end{frame}

% ==========================================
% Section 4: Experimental Environment & Baselines
% ==========================================
% \begin{frame}{Experimental Setup}
%     \begin{columns}[T]
%         \begin{column}{0.48\textwidth}
%             \small
%             \begin{itemize}
%                 \item Five standard vision benchmarks are selected: MNIST, FMNIST, CIFAR-10, CIFAR-100, and FEMNIST.
%                 \item Data distributions across clients are skewed using a Dirichlet concentration parameter $\alpha \in \{0.1, 1, 10\}$.
%                 \item Bandwidth bottlenecks are simulated using a uniform distribution $U(5, 20)$~Mbps.
%             \end{itemize}
%         \end{column}
%         \begin{column}{0.48\textwidth}
%             \centering
%             \includegraphics[width=\textwidth,height=0.65\textheight,keepaspectratio]{figures/jimenez_fig7.png}
%             \vspace{0.1cm}

%             {\fontsize{5.5pt}{7pt}\selectfont Jimenez et al. (2024), Figure 7: Taxonomy of different types of skew and heterogeneity.\\}
%         \end{column}
%     \end{columns}
%     \vfill
%     \centering
%     {\fontsize{8pt}{9.5pt}\selectfont Statistical and system heterogeneity are modeled simultaneously to reflect realistic edge environments.}
% \end{frame}

\begin{frame}{Classification of Communication-Efficient FL Related Works}
    \begin{columns}[T]
        \begin{column}{0.48\textwidth}
            \begin{block}{Foundational Baselines}
                \small FedAvg, FedProx
            \end{block}
            \vspace{0.2cm}
            \begin{block}{Static Compression}
                \small DAdaQuant, FedPAQ
            \end{block}
            \vspace{0.2cm}
            \begin{block}{Pure Distillation}
                \small FedKD, FedDistill
            \end{block}
        \end{column}
        \begin{column}{0.48\textwidth}
            \begin{block}{Adaptive Compression}
                \small AdaGQ, LAQ-HC
            \end{block}
            \vspace{0.2cm}
            \begin{block}{Direct SOTA Competitors}
                \small DynFed, AdaDQ-KD, FedDT
            \end{block}
        \end{column}
    \end{columns}
    \vfill
    \centering
    {\fontsize{8pt}{9.5pt}\selectfont An 11-algorithm evaluation stack benchmarks FedMAQ against compression, distillation, and adaptive SOTA.}
\end{frame}

% ==========================================
% Section 5: Progress & Next Steps
% ==========================================
\begin{frame}{Manuscript Progress}
    \centering
    \vfill
    \begin{table}
        \fontsize{8pt}{11pt}\selectfont
        \begin{tabular}{lp{8.0cm}}
            \toprule
            \textbf{Chapter}                          & \textbf{Status \& Accomplishments} \\
            \midrule
            \textbf{Chapter 1: Introduction}          & Outlined                           \\
            \textbf{Chapter 2: Related Works}         & Outlined                           \\
            \textbf{Chapter 3: Theoretical Framework} & Outlined                           \\
            \textbf{Chapter 4: Methodology}           & Outlined                           \\
            \bottomrule
        \end{tabular}
    \end{table}
    \vfill
\end{frame}

\begin{frame}{Next Meeting Commitments}
    \begin{itemize}
        \item Configure the Flower-based Federated Learning simulation environment for Phase 1 baselines.
        \item Implement Dirichlet non-IID partitioning pipelines across the selected image classification datasets.
    \end{itemize}
\end{frame}

% Slide 8: Q&A / Standout Frame
\begin{frame}[standout]
    Thank You!
    \vspace{0.5cm}

    \begin{center}
        \large Questions \& Answers
    \end{center}
\end{frame}

\end{document}
